\section{Initial Situation}

\lead{As of 28th November 1975, the international community has become increasingly aware of the widespread political repression taking place across the Southern Cone.}

Military governments have intensified campaigns against political opponents, resulting in a growing number of reports of arbitrary detentions, enforced disappearances, torture, and political persecution. Nevertheless, these developments are generally perceived as isolated domestic matters rather than evidence of a coordinated regional effort.

Human rights organizations, members of the press, and political exiles continue to raise concerns regarding these abuses, particularly in Chile, where international scrutiny has steadily increased since the 1973 coup d’état. The activities of the DINA and the large-scale exile of Chilean dissidents are well known abroad, reinforcing concerns over the country’s human rights record.

The political climate across Latin America has become increasingly hostile to those who oppose the region's military governments. The continued expansion of authoritarian rule, coupled with the systematic persecution of political dissidents, has profoundly altered the balance of power, forcing opposition movements to reconsider both their objectives and the means through which they pursue them. As governments strengthen their control over state institutions, the space for political resistance continues to narrow.

Against this backdrop, those committed to defending democratic principles, social justice, and anti-capitalist ideals face an increasingly uncertain future. The preservation of these values has become inseparable from the broader struggle against authoritarianism, requiring opposition leaders to evaluate new forms of cooperation, reassess existing strategies, and determine how best to respond to a rapidly changing political environment. For many, the question is no longer whether resistance is necessary, but how it can be carried out and survive.

Recent reports indicate that a meeting of several Latin American leaders has taken place, although little is known about the substance of the discussions or the agreements that may have been reached. The possibility of closer cooperation between regional governments has generated growing concern among opposition movements, reinforcing the perception that the challenge they face is no longer confined to individual states but may increasingly extend beyond national borders. In this context, safeguarding both their political cause and their own survival has become an essential priority, with many believing that extraordinary circumstances may require equally extraordinary measures.
